\documentclass[11pt,a4paper]{article}
\usepackage{gtu-el-yazisi}
\begin{document}
\notbilgi{Sayısal Haberleşme}{Olasılık teorisi, rastgele değişkenler ve prosesler}{Ders2.pdf}

\sayfa{1}
\section{Olasılık teorisi, Rastgele Değişkenler ve Rastgele Prosesler}

\vurgu[turuncu]{Deney (experiement)}: Zar atışı

\vurgu[turuncu]{Sonuç (outcome)}

\vurgu[turuncu]{Örnek uzayı (sample space)} : $S=\{1,2,3,4,5,6\}$

\vurgu[turuncu]{Olay (event)} : $A=\{\text{Çift sayılar}\}=\{2,4,6\}$

$A\subset S$

\elcizimi[0.45]{sekiller/Ders2-s1-1.png}[1]

\begin{align*}
P(\overset{\text{olay}}{A}) &= \sum_{\forall s_i\in A} P(s_i) = P(2)+P(4)+P(6)\\
&= \frac16+\frac16+\frac16 = \frac36 = \frac12
\end{align*}

Eşit olasılıklı sonuçlar için
\[ P(A)=\frac{|A|}{|S|}=\frac{\text{A'nın eleman sayısı}}{\text{S'nin \ '' \ ''}}=\frac{3}{6}=\frac{1}{2} \]

Bir olayın olması için olaydaki tek bir sonucun olması yeterli.

Kesin olay (Certain event) $=S \rightarrow P(S)=1$

İmkansız olay (Impossible event) $=\emptyset \rightarrow P(\emptyset)=0$

\elcizimi[0.3]{sekiller/Ders2-s1-2.png}[1]

\begin{align*}
&P(A\cup B)=P(A)+P(B)-P(A\cap B)\\
&A\cap B=\emptyset \Rightarrow A\ \text{ve}\ B\ \text{kesişimsiz olaylar}\\
&\qquad\text{(mutually exclusive, disjoint)}
\end{align*}

\begin{align*}
A\cap B=\emptyset \Rightarrow P(A\cup B) &= P(A)+P(B)-P(\underbrace{A\cap B}_{0})\\
P(A\cup B) &= P(A)+P(B)
\end{align*}

$A'=\{a \mid a\in S,\ a\notin A\}$ \quad uzayın A dışındaki tüm elemanlar
\[ A\cup A'=S \]
\[ A\cap A'=\emptyset \]

\sayfa{2}
\begin{align*}
S &= A\cup A'\\
P(S) &= P(A\cup A') = P(A)+P(A')-P(\underbrace{A\cap A'}_{\emptyset})\\
1 &= P(A)+P(A')\\
\Rightarrow\ P(A') &= 1-P(A)
\end{align*}

\[ A \overset{\text{Bağımsızsa (independent)}}{\perp\!\!\!\perp} B \iff P(A\cap B)=P(A)\cdot P(B) \]

\section{Rastgele Değişken (RD)}
$X(s)$ : RD

\elcizimi[0.9]{sekiller/Ders2-s2-1.png}[2]

\[ \overset{\text{Rastgele değişken}}{X}(s)=x \]
$x$ : $X$ RD'nin belli bir değeri.

örnek $Y$, $y$

\section{Olasılık Dağılım Fonksiyonu (ODF)}
Cumulative Distribution function - CDF
\[ F_X(x)=\overset{\text{olasılık}}{P}(X\le x) \]

\elcizimi[0.5]{sekiller/Ders2-s2-2.png}[2]

\begin{itemize}
\item $F_X(-\infty)=0$
\item $F_X(+\infty)=1$
\item $0\le F_X(x)\le 1$
\item $x_1\le x_2 \Rightarrow F_X(x_1)\le F_X(x_2)$ \hfill Azalmayan fonksiyon (Nondecreasing function)
\end{itemize}

\section{Olasılık Yoğunluk Fonk. (OYF)}
Probability Density Function (PDF)

Bir rastgele değişkenin belli bir bölgede olmasının olasılığını, altında kalan alanla gösteren fonksiyona denir.

$f_X(x)$ OYF

\sayfa{3}
\begin{align*}
P(X\in A) &= \int_{A} f_X(x)\,dx\\
P(X\le x) &= \int_{-\infty}^{x} \overset{\text{OYF}}{f_X(t)}\,dt = \overset{\text{ODF}}{F_X(x)}
\end{align*}

\[ f_X(x)=\frac{d}{dx}F_X(x) \]

\begin{itemize}
\item $\displaystyle\int_{-\infty}^{\infty} f_X(x)\,dx=1$
\item $\forall x\quad f_X(x)\ge 0$ \hfill azalmayan fonksiyonun türevi pozitiftir
\end{itemize}

\[ \text{1)}\quad P(X\le a)=\int_{-\infty}^{a} f_X(u)\,du=F_X(a) \]

\elcizimi[0.5]{sekiller/Ders2-s3-1.png}[3]

\begin{align*}
\text{2)}\quad P(\underbrace{X>a}_{A}) &= \int_{a}^{\infty} f_X(u)\,du = 1-\underbrace{\int_{-\infty}^{a} f_X(u)\,du}_{P(A')}\\
&= 1-F_X(a)
\end{align*}

\elcizimi[0.5]{sekiller/Ders2-s3-2.png}[3]

\begin{align*}
\text{3)}\quad P(a<X\le b) &= \int_{a}^{b} f_X(u)\,du = \int_{-\infty}^{b} f_X(u)\,du-\int_{-\infty}^{a} f_X(u)\,du\\
&= F_X(b)-F_X(a)
\end{align*}

\elcizimi[0.5]{sekiller/Ders2-s3-3.png}[3]

\sayfa{4}
\section{Düzgün Dağılım (Uniform distribution)}

\[ f_X(x)=\begin{cases} \dfrac{1}{b-a}, & a<x<b\\[4pt] 0, & \text{diğer} \end{cases} \]

\elcizimi[0.5]{sekiller/Ders2-s4-1.png}[4]

\[ F_X(x)=\int_{-\infty}^{x} f_X(u)\,du \]

\[ \textcolor{kirmizikalem}{\underline{\text{1. bölge}}}\ :\ x<a\footnotemark \Rightarrow f_X(x)=0 \Rightarrow F_X(x)=0 \]
\footnotetext{Claude notu: orijinalde siyah kalemle $x<0$ yazılmış, turuncu kalemle $x<a$ olarak düzeltilmiş.}

\begin{align*}
\textcolor{kirmizikalem}{\underline{\text{2. bölge}}}\ :\ a<x<b \Rightarrow F_X(x) &= \int_{-\infty}^{x} f_X(u)\,du = \int_{a}^{x}\frac{1}{b-a}\,du\\
&= \left.\frac{u}{b-a}\right|_{a}^{x}\\
&= \frac{x-a}{b-a}
\end{align*}

\begin{align*}
\textcolor{kirmizikalem}{\underline{\text{3. bölge}}}\ :\ x>b \Rightarrow F_X(x) &= \int_{-\infty}^{x} f_X(u)\,du = \int_{a}^{b}\frac{1}{b-a}\,du\\
&= \left.\frac{u}{b-a}\right|_{a}^{b} = \frac{b-a}{b-a}=1
\end{align*}

\[ F_X(x)=\begin{cases} 0, & x<a\\[2pt] \dfrac{x-a}{b-a}, & a<x<b\\[4pt] 1, & x\ge b \end{cases} \]

\elcizimi[0.5]{sekiller/Ders2-s4-2.png}[4]

\section{Normal Dağılım (Gauss)}

\[ f_X(x)=\frac{1}{\sqrt{2\pi}\,\sigma_X}\,e^{-\frac{(x-\mu_X)^2}{2\sigma_X^2}} = N(\mu_X,\sigma_X^2) \]

$\mu_X$ : ortalama değer (mean)

$\sigma_X$ : Standart sapma

$\sigma_X^2$ = Varyans

\sayfa{5}
\elcizimi[0.55]{sekiller/Ders2-s5-1.png}[5]

\[ F_X(x)=\int_{-\infty}^{x} f_X(u)\,du=\frac{1}{\sqrt{2\pi}\,\sigma_X}\int_{-\infty}^{x} e^{-\frac{(u-\mu_X)^2}{2\sigma_X^2}}\,du \]

\section{Standart Gauss Dağılımı}

\[ N(0,1):\quad f_X(x)=\frac{1}{\sqrt{2\pi}}\,e^{-\frac{x^2}{2}} \]

\elcizimi[0.5]{sekiller/Ders2-s5-2.png}[5]

\[ Q(x)=P(X>x)=\int_{x}^{\infty} f_X(t)\,dt=\frac{1}{\sqrt{2\pi}}\int_{x}^{\infty} e^{-\frac{t^2}{2}}\,dt\ ,\quad x\ge 0 \]

S 104\okunamadi{5} Table B1 $x\ge 0$ \hfill Q tablosu

\[ Q(0)=0.5 \]
\[ Q(-x)=1-Q(x) \]

\[ \frac{u-\mu_X}{\sigma_X}=t \]
\[ u=t\sigma_X+\mu_X \qquad du=\sigma_X\,dt \]
\begin{align*}
u=-\infty &\rightarrow t=-\infty\\
u=x &\rightarrow t=\frac{x-\mu_X}{\sigma_X}
\end{align*}

\begin{align*}
F_X(x) &= \int_{-\infty}^{x} f_X(u)\,du = \frac{1}{\sqrt{2\pi}\,\sigma_X}\int_{-\infty}^{x} e^{-\frac{(u-\mu_X)^2}{2\sigma_X^2}}\,du\\
&= \frac{1}{\sqrt{2\pi}\,\sigma_X}\int_{-\infty}^{\frac{x-\mu_X}{\sigma_X}} e^{-\frac{t^2}{2}}\,\sigma_X\,dt\\
&= \kutula{1-\frac{1}{\sqrt{2\pi}}\underbrace{\int_{\frac{x-\mu_X}{\sigma_X}}^{\infty} e^{-t^2/2}\,dt}_{Q\left(\frac{x-\mu_X}{\sigma_X}\right)}} = 1-Q\left(\frac{x-\mu_X}{\sigma_X}\right)
\end{align*}

\sayfa{6}
\section{Beklenen Değer (Expectation)}

\[ \mu_X=E[X]=\int_{-\infty}^{+\infty} x\,f_X(x)\,dx \]
$\hookrightarrow$ X'in beklenen değeri (ortalaması)

\[ \underbrace{E[X^n]}_{\text{X'in n. momenti}}=\int_{-\infty}^{+\infty} x^n f_X(x)\,dx \]

\[ E[g(x)]=\int_{-\infty}^{+\infty} g(x)\,f_X(x)\,dx \]

\begin{align*}
&g(X)=(X-\mu_X)^2 \Rightarrow \sigma_X^2=E\left[(X-\mu_X)^2\right]=\int_{-\infty}^{\infty}(x-\mu_X)^2 f_X(x)\,dx\\
&\qquad\qquad = \underbrace{\int x^2 f_X(x)\,dx}_{=\,E[X^2]} + \mu_X^2\underbrace{\int f_X(x)\,dx}_{1} - 2\mu_X\underbrace{\int x f_X(x)\,dx}_{\mu_X}
\end{align*}

\[ \sigma_X^2=E[X^2]-\mu_X^2=E[X^2]-(E[X])^2 \]

\section{karekteristik fonk}

\[ \phi_X(\omega)=E\left[e^{j\omega X}\right]=\int_{-\infty}^{+\infty} e^{j\omega x} f_X(x)\,dx \]

\[ \kutula{E[X^n]=(-j)^n\left.\frac{d^n\phi_X(\omega)}{d\omega^n}\right|_{\omega=0}} \]

$e^{\omega x}$ olduğunda moment üreten fonksiyon deniyor

Son zamanlarda

Moment üreten fonk (MÜF):
\[ M_X(s)=E\left[e^{sX}\right]=\int_{-\infty}^{\infty} e^{sx} f_X(x)\,dx \qquad \searrow\ \text{laplace dönüşümü} \]

\[ E[X^n]=\left.\frac{d^n}{ds^n}M_X(s)\right|_{s=0} \]

\textcolor{kirmizikalem}{Örnek:} Düzgün dağılım

\elcizimi[0.45]{sekiller/Ders2-s6-1.png}[6]

\sayfa{7}
\begin{align*}
\mu_X=E[X] &= \int_{-\infty}^{\infty} x f_X(x)\,dx\\
&= \int_{a}^{b} x\,\frac{1}{b-a}\,dx = \frac{1}{b-a}\left.\frac{x^2}{2}\right|_{a}^{b}\\
&= \frac{b^2-a^2}{2(b-a)} = \frac{a+b}{2}
\end{align*}

\begin{align*}
E[X^2] &= \int_{-\infty}^{\infty} x^2 f_X(x)\,dx = \int_{a}^{b} x^2\,\frac{1}{b-a}\,dx = \frac{1}{b-a}\left.\frac{x^3}{3}\right|_{a}^{b}\\
&= \frac{b^3-a^3}{3(b-a)} = \frac{a^2+ab+b^2}{3}
\end{align*}

\begin{align*}
\sigma_X^2 &= E[X^2]-\mu_X^2 = \frac{a^2+ab+b^2}{3}-\left(\frac{a+b}{2}\right)^2\\
&= \underset{(4)}{\frac{a^2+ab+b^2}{3}} - \underset{(3)}{\frac{a^2+b^2+2ab}{4}}\\
&= \frac{a^2-2ab+b^2}{12} = \frac{(b-a)^2}{12}
\end{align*}

\formulkutu[önemli]{\begin{aligned}\mu_X&=\frac{a+b}{2}\\ \sigma_X^2&=\frac{(b-a)^2}{12}\end{aligned}}

\section{Rastgele Prosesler (Random, stochastic)}

\elcizimi[0.9]{sekiller/Ders2-s7-1.png}[7]

RP : $X(s,t)=X(t)$

$\hookrightarrow$ örnek fonksiyonlar topluluğu.

$X(t_k)=X_{t_k}=X_k$ \quad RD

\sayfa{8}
\[ E[X_{t_k}]=\int_{-\infty}^{\infty} x_{t_k}\cdot f_{X_{t_k}}(x_{t_k})\,dx_{t_k}=\mu_{X_{t_k}}\qquad\text{RP'in beklenen değeri.} \]

RP'in özilişki (auto correlation) fonk

\begin{align*}
R(t_1,t_2) &= E[\underbrace{X(t_1)X(t_2)}_{x_{t_1}\cdot x_{t_2}}]\\
&= \int_{-\infty}^{\infty}\!\int_{-\infty}^{\infty} x_{t_1}\cdot x_{t_2}\,\underbrace{f_X(x_{t_1},x_{t_2})}_{\substack{x_{t_1}\ \text{ve}\ x_{t_2}\\ \text{RP'lerinin ortak OYF}\\ \text{(joint PDF)}}}\,dx_{t_1}\,dx_{t_2}
\end{align*}

\section{Stasyonerlik}

\textcolor{kirmizikalem}{-- Sıkı Stasyoner (SS)} \quad : Bütün momentler zamana bağlı değil

Strict sense stationary (SSS)

\textcolor{kirmizikalem}{-- Gevşek Stasyoner (GS)} \quad : ilk 2 moment zamana bağlı değil

Wide sense stationary (WSS)

Bir RP aşağıdaki şartları sağlıyorsa GS'dir

\[ \text{1)}\quad E[X(t)]=\mu_X=\text{sabit} \]
\[ \text{2)}\quad R_X(t_1,t_2)=R_X(\underbrace{t_1-t_2}_{\tau})=R_X(\tau)=E[X(t)\,X(t-\tau)] \]
2.moment $\tau=t_1-t_2$ ye bağlı.

Process \qquad SS $\Rightarrow$ GS

\qquad\qquad\qquad\quad $\not\Leftarrow$

\qquad\qquad\qquad\quad $\Leftarrow$ Gauss RP sağlar

Haberleşme sistemlerinde GS önemli:

\section{Özellikler}

\begin{itemize}
\item $R_X(-\tau)=R_X(\tau)$ \quad çift fonk.
\item $R_X(0)\ge|R_X(\tau)|$ $\hookrightarrow$ complex sayılar için
\item $R_X(0)=E\left[|X(t)|^2\right]$ \quad Güç
\end{itemize}

Rastgele Processler Güç işaretleridir burda istatistiksel ortalama alınıyor

\sayfa{9}
\begin{itemize}
\item $R_X(\tau)\overset{F}{\longleftrightarrow}S_X(f)$ \quad GYF \quad güç yoğunluk fonksiyonu
\end{itemize}

\section{GYF}

RP Güç işareti $\rightarrow$ GYF, $S_X(f)$

\begin{itemize}
\item $R_X(\tau)\overset{F}{\longleftrightarrow}S_X(f)$ \quad $\dfrac{\text{Watt}}{\text{Hertz}}$
\item $S_X(f)\ge 0$ \quad (Reel değerli)
\end{itemize}

\begin{itemize}
\item Prosesin ortalama gücü
\[ P_X=R_X(0)=\int_{-\infty}^{+\infty} S_X(f)\,df \]
\item Reel değerli $X(t)$ RP için $S_X(-f)=S_X(f)$
\end{itemize}

\section{Beyaz Gürültü RP}
(White Noise RP)

\[ n(t)\rightarrow S_n(f)=\frac{N_0}{2}=\text{sabit}\qquad \frac{\text{Watt}}{\text{Hz}}\quad \forall f \]

\elcizimi[0.85]{sekiller/Ders2-s9-1.png}[9]

\[ R_n(\tau)=\frac{N_0}{2}\,\delta(\tau) \]

\elcizimi[0.6]{sekiller/Ders2-s9-2.png}[9]

Gücü sonsuz

Gerçekte gücü sonsuz güce sahip işaret yok

Bütün frekanslarda olduğu için Beyaz gürültü denir.

(Prizma)

\sayfa{10}
\elcizimi[0.9]{sekiller/Ders2-s10-1.png}[10]

\[ \underbrace{r(t_k)}_{\substack{\text{Gauss}\\ \text{RD}}} = \underbrace{s(t_k)}_{\substack{a_i\\ \text{Deterministic}}} + \underbrace{n(t_k)}_{n} \ \rightarrow\ \text{Gauss R.D} \]

\[ \underset{\substack{\downarrow\\ N(a,\sigma_n^2)}}{z}=a+n\ \rightarrow\ N(0,\sigma_n^2) \]

\[ f_n(n)=\frac{1}{\sqrt{2\pi}\,\sigma_n}\,e^{-\frac{n^2}{2\sigma_n^2}} \]

\elcizimi[0.3]{sekiller/Ders2-s10-2.png}[10]

\[ f_z(z)=\frac{1}{\sqrt{2\pi}\,\sigma_n}\,e^{-\frac{(z-a)^2}{2\sigma_n^2}} \]

\elcizimi[0.45]{sekiller/Ders2-s10-3.png}[10]

Alıcı alırken $z$ pozitif mi negatif mi diye bakıyor

$s(t)$ pozitifken \quad $z(t)$'nin negatif olma olasılığını

'' \quad negatifken \quad '' \quad pozitif \quad '' \quad ''

bulup toplayıp ikiye bölerek yapıyor

\section{LZD sistemler}

\elcizimi[0.35]{sekiller/Ders2-s10-4.png}[10]

\begin{align*}
y(t) &= x(t)\otimes h(t)\\
&= \int_{-\infty}^{\infty} x(\tau)\,h(t-\tau)\,d\tau
\end{align*}

\begin{tabular}{@{}p{0.4\textwidth}l@{}}
$X(t)$ RP & $Y(t)$ RP\\[4pt]
$S_X(f)$ & $S_Y(f)=S_X(f)\,|H(f)|^2$
\end{tabular}

\section{İşaretler ve Sistemler}

$(-T/2,\ T/2)$ arasında harcanan enerji
\[ E_X^T=\int_{-T/2}^{T/2}|x(t)|^2\,dt \]

\sayfa{11}
güç \quad $P_X^T=\dfrac{E_X^T}{T}$

\[ E_X=\lim_{T\to\infty} E_X^T=\lim_{T\to\infty}\int_{-T/2}^{T/2}|x(t)|^2\,dt=\int_{-\infty}^{\infty}|x(t)|^2\,dt \]

\[ P_X=\lim_{T\to\infty}\frac{1}{T}\int_{-T/2}^{T/2}|x(t)|^2\,dt \]

\begin{align*}
E_X<\infty &\Rightarrow x(t)\ \text{enerji işareti}\\
P_X<\infty &\Rightarrow \text{''}\ \text{güç işareti}
\end{align*}

\[ x(t)\ \text{periyodik}\ \Rightarrow\ P_X=\frac{1}{T_0}\int_{T_0}|x(t)|^2\,dt \]

\[ x(t)=\begin{cases}A\cos(2\pi f_0 t)\\ A\sin(2\pi f_0 t)\end{cases}\Rightarrow\ P_X=\frac{A^2}{2}\ \text{watt} \]

\end{document}
